\documentclass{report}
\usepackage{amsmath,amssymb}
\setlength{\parindent}{0mm}
\setlength{\parskip}{1em}
\begin{document}
\begin{center}
\rule{6in}{1pt} \
{\large Jeffery Allen \\
{\bf Fluidity Based Formulation of Nonlinear Stokes Flow for Ice Sheets using FOSLS}}

2210 Walnut Apt 1 \\ Boulder \\ CO 80302
\\
{\tt jeffery.allen@colorado.edu}\\
Thomas Manteuffel \\
Harihar Rajaram\end{center}

This talk describes two First-order System Least Squares (FOSLS)
formulations of a nonlinear Stokes flow model for glaciers and ice
sheets. In Glen's law, the most commonly used constitutive equation for
ice rheology, the ice viscosity becomes infinite as the velocity
gradients (strain rates) approach zero, which typically occurs near the
ice surface where deformation rates are low, or when the basal slip
velocities are high. The computational difficulties associated with the
infinite viscosity are often overcome by an arbitrary modification of
Glen's law that bounds the maximum viscosity. In this talk, two FOSLS
formulations (the viscosity formulation and fluidity formulation) are
presented. The viscosity formulation is a FOSLS representation of the
standard nonlinear Stokes problem. The new fluidity formulation exploits
the fact that only the product of the viscosity and strain rate appears
in the nonlinear Stokes problem, a quantity that, in fact, approaches
zero as the strain rate goes to zero. The fluidity formulation is
expressed in terms of a new set of variables and overcomes the problem of
infinite viscosity. The new formulation is well posed and the
linearization is essentially $H^1$-elliptic around solutions for which
the viscosity is bounded. A Nested Iteration (NI) Newton-FOSLS-AMG
approach is used to solve both nonlinear Stokes problems, in which most
of the iterations are performed on the coarsest grid. Both formulations
demonstrate optimal finite element convergence. However, the fluidity
formulation is more accurate. The fluidity formulation involves linear
systems that are more amenable to solution by AMG. Further improvement in
computational cost is achieved using local adaptive refinement.


\end{document}
